% F14：本书原创矢量示意；依据所在节的定义、证明或明确模型。
\begin{figure}[htbp]
\centering
\begin{tikzpicture}[x=1cm,y=1cm,>=stealth,every node/.style={font=\small},line width=.65pt]
\node[] at (3.0,3.4) {缩放参数 \(r=1\)};
\draw[AnalysisBlue,thick] plot coordinates {(1.0000,2.6000) (1.0506,2.5500) (1.1013,2.5013) (1.1519,2.4539) (1.2025,2.4077) (1.2532,2.3629) (1.3038,2.3193) (1.3544,2.2770) (1.4051,2.2360) (1.4557,2.1962) (1.5063,2.1578) (1.5570,2.1206) (1.6076,2.0847) (1.6582,2.0501) (1.7089,2.0168) (1.7595,1.9847) (1.8101,1.9539) (1.8608,1.9245) (1.9114,1.8963) (1.9620,1.8693) (2.0127,1.8437) (2.0633,1.8194) (2.1139,1.7963) (2.1646,1.7745) (2.2152,1.7540) (2.2658,1.7348) (2.3165,1.7168) (2.3671,1.7001) (2.4177,1.6848) (2.4684,1.6707) (2.5190,1.6578) (2.5696,1.6463) (2.6203,1.6361) (2.6709,1.6271) (2.7215,1.6194) (2.7722,1.6130) (2.8228,1.6079) (2.8734,1.6040) (2.9241,1.6014) (2.9747,1.6002) (3.0253,1.6002) (3.0759,1.6014) (3.1266,1.6040) (3.1772,1.6079) (3.2278,1.6130) (3.2785,1.6194) (3.3291,1.6271) (3.3797,1.6361) (3.4304,1.6463) (3.4810,1.6578) (3.5316,1.6707) (3.5823,1.6848) (3.6329,1.7001) (3.6835,1.7168) (3.7342,1.7348) (3.7848,1.7540) (3.8354,1.7745) (3.8861,1.7963) (3.9367,1.8194) (3.9873,1.8437) (4.0380,1.8693) (4.0886,1.8963) (4.1392,1.9245) (4.1899,1.9539) (4.2405,1.9847) (4.2911,2.0168) (4.3418,2.0501) (4.3924,2.0847) (4.4430,2.1206) (4.4937,2.1578) (4.5443,2.1962) (4.5949,2.2360) (4.6456,2.2770) (4.6962,2.3193) (4.7468,2.3629) (4.7975,2.4077) (4.8481,2.4539) (4.8987,2.5013) (4.9494,2.5500) (5.0000,2.6000)};
\draw[AnalysisBlue,thick] plot coordinates {(1.0000,0.6000) (1.0506,0.6500) (1.1013,0.6987) (1.1519,0.7461) (1.2025,0.7923) (1.2532,0.8371) (1.3038,0.8807) (1.3544,0.9230) (1.4051,0.9640) (1.4557,1.0038) (1.5063,1.0422) (1.5570,1.0794) (1.6076,1.1153) (1.6582,1.1499) (1.7089,1.1832) (1.7595,1.2153) (1.8101,1.2461) (1.8608,1.2755) (1.9114,1.3037) (1.9620,1.3307) (2.0127,1.3563) (2.0633,1.3806) (2.1139,1.4037) (2.1646,1.4255) (2.2152,1.4460) (2.2658,1.4652) (2.3165,1.4832) (2.3671,1.4999) (2.4177,1.5152) (2.4684,1.5293) (2.5190,1.5422) (2.5696,1.5537) (2.6203,1.5639) (2.6709,1.5729) (2.7215,1.5806) (2.7722,1.5870) (2.8228,1.5921) (2.8734,1.5960) (2.9241,1.5986) (2.9747,1.5998) (3.0253,1.5998) (3.0759,1.5986) (3.1266,1.5960) (3.1772,1.5921) (3.2278,1.5870) (3.2785,1.5806) (3.3291,1.5729) (3.3797,1.5639) (3.4304,1.5537) (3.4810,1.5422) (3.5316,1.5293) (3.5823,1.5152) (3.6329,1.4999) (3.6835,1.4832) (3.7342,1.4652) (3.7848,1.4460) (3.8354,1.4255) (3.8861,1.4037) (3.9367,1.3806) (3.9873,1.3563) (4.0380,1.3307) (4.0886,1.3037) (4.1392,1.2755) (4.1899,1.2461) (4.2405,1.2153) (4.2911,1.1832) (4.3418,1.1499) (4.3924,1.1153) (4.4430,1.0794) (4.4937,1.0422) (4.5443,1.0038) (4.5949,0.9640) (4.6456,0.9230) (4.6962,0.8807) (4.7468,0.8371) (4.7975,0.7923) (4.8481,0.7461) (4.8987,0.6987) (4.9494,0.6500) (5.0000,0.6000)};
\draw[orange!80!black,dashed] plot coordinates {(0.7000,1.6000) (5.3000,1.6000)};
\node[] at (3.0,-0.35) {\(y=\pm r x^2\)};
\node[] at (9.2,3.4) {缩放参数 \(r=0.25\)};
\draw[AnalysisBlue,thick] plot coordinates {(7.2000,1.8500) (7.2506,1.8375) (7.3013,1.8253) (7.3519,1.8135) (7.4025,1.8019) (7.4532,1.7907) (7.5038,1.7798) (7.5544,1.7692) (7.6051,1.7590) (7.6557,1.7491) (7.7063,1.7394) (7.7570,1.7301) (7.8076,1.7212) (7.8582,1.7125) (7.9089,1.7042) (7.9595,1.6962) (8.0101,1.6885) (8.0608,1.6811) (8.1114,1.6741) (8.1620,1.6673) (8.2127,1.6609) (8.2633,1.6548) (8.3139,1.6491) (8.3646,1.6436) (8.4152,1.6385) (8.4658,1.6337) (8.5165,1.6292) (8.5671,1.6250) (8.6177,1.6212) (8.6684,1.6177) (8.7190,1.6145) (8.7696,1.6116) (8.8203,1.6090) (8.8709,1.6068) (8.9215,1.6048) (8.9722,1.6032) (9.0228,1.6020) (9.0734,1.6010) (9.1241,1.6004) (9.1747,1.6000) (9.2253,1.6000) (9.2759,1.6004) (9.3266,1.6010) (9.3772,1.6020) (9.4278,1.6032) (9.4785,1.6048) (9.5291,1.6068) (9.5797,1.6090) (9.6304,1.6116) (9.6810,1.6145) (9.7316,1.6177) (9.7823,1.6212) (9.8329,1.6250) (9.8835,1.6292) (9.9342,1.6337) (9.9848,1.6385) (10.0354,1.6436) (10.0861,1.6491) (10.1367,1.6548) (10.1873,1.6609) (10.2380,1.6673) (10.2886,1.6741) (10.3392,1.6811) (10.3899,1.6885) (10.4405,1.6962) (10.4911,1.7042) (10.5418,1.7125) (10.5924,1.7212) (10.6430,1.7301) (10.6937,1.7394) (10.7443,1.7491) (10.7949,1.7590) (10.8456,1.7692) (10.8962,1.7798) (10.9468,1.7907) (10.9975,1.8019) (11.0481,1.8135) (11.0987,1.8253) (11.1494,1.8375) (11.2000,1.8500)};
\draw[AnalysisBlue,thick] plot coordinates {(7.2000,1.3500) (7.2506,1.3625) (7.3013,1.3747) (7.3519,1.3865) (7.4025,1.3981) (7.4532,1.4093) (7.5038,1.4202) (7.5544,1.4308) (7.6051,1.4410) (7.6557,1.4509) (7.7063,1.4606) (7.7570,1.4699) (7.8076,1.4788) (7.8582,1.4875) (7.9089,1.4958) (7.9595,1.5038) (8.0101,1.5115) (8.0608,1.5189) (8.1114,1.5259) (8.1620,1.5327) (8.2127,1.5391) (8.2633,1.5452) (8.3139,1.5509) (8.3646,1.5564) (8.4152,1.5615) (8.4658,1.5663) (8.5165,1.5708) (8.5671,1.5750) (8.6177,1.5788) (8.6684,1.5823) (8.7190,1.5855) (8.7696,1.5884) (8.8203,1.5910) (8.8709,1.5932) (8.9215,1.5952) (8.9722,1.5968) (9.0228,1.5980) (9.0734,1.5990) (9.1241,1.5996) (9.1747,1.6000) (9.2253,1.6000) (9.2759,1.5996) (9.3266,1.5990) (9.3772,1.5980) (9.4278,1.5968) (9.4785,1.5952) (9.5291,1.5932) (9.5797,1.5910) (9.6304,1.5884) (9.6810,1.5855) (9.7316,1.5823) (9.7823,1.5788) (9.8329,1.5750) (9.8835,1.5708) (9.9342,1.5663) (9.9848,1.5615) (10.0354,1.5564) (10.0861,1.5509) (10.1367,1.5452) (10.1873,1.5391) (10.2380,1.5327) (10.2886,1.5259) (10.3392,1.5189) (10.3899,1.5115) (10.4405,1.5038) (10.4911,1.4958) (10.5418,1.4875) (10.5924,1.4788) (10.6430,1.4699) (10.6937,1.4606) (10.7443,1.4509) (10.7949,1.4410) (10.8456,1.4308) (10.8962,1.4202) (10.9468,1.4093) (10.9975,1.3981) (11.0481,1.3865) (11.0987,1.3747) (11.1494,1.3625) (11.2000,1.3500)};
\draw[orange!80!black,dashed] plot coordinates {(6.9000,1.6000) (11.5000,1.6000)};
\node[] at (9.2,-0.35) {\(y=\pm r x^2\)};

\end{tikzpicture}
\caption[相切曲线的同一切线与重数二]{相切曲线的同一切线与重数二。两条曲线在原点有同一近似切线，但每支贡献一份长度，密度为二，缩放极限带重数二。这不反驳几乎处处的单位密度结论。}
\label{b1:fig:visual-f14}
\end{figure}
